% !TEX root = relatorio_oncologia.tex
\documentclass[11pt,a4paper]{article}
\usepackage[utf8]{inputenc}
\usepackage[portuguese]{babel}
\usepackage{graphicx}
\usepackage{caption}
\usepackage{geometry}
\usepackage{amsmath}
\geometry{a4paper, margin=2.1cm}

\begin{document}

\begin{titlepage}
    \centering
    \vspace*{1cm}
    
    {\large universidade de aveiro \par}
    {\large departamento de electrónica, \par deti telecomunicações e informática \par}
    
    \vspace{2.5cm}
    
    {\Large Métodos Probabilísticos para Engenharia Informática \par 2025/26 \par}
    
    \vspace{2.5cm}
    
    {\huge\bfseries Projeto Final \par Sistema de Apoio à Decisão Clínico-Oncológica \par através de Algoritmos Probabilísticos \par}
    
    \vspace{3cm}
    
    
    {\Large Filipe Nogueira - 119495 \par João Tomásio - 120132 \par}
    
    \vfill
    
    {\large Professor Regente: \par António Joaquim da Silva Teixeira (ajst@ua.pt) \par}
    
    \vspace*{1cm}
\end{titlepage}

\newpage

\section*{Introdução}
Este relatório descreve a implementação de um sistema de apoio à decisão Clínico‑Oncológica desenvolvido em MatLab, com o objetivo de integrar técnicas probabilísticas para a análise de dados genómicos. Utilizando Filtros bloom para triagem eficiente, MinHash para análise de similaridade e o classificador Naïve Bayes para prever o impacto clínico de mutações, oferecendo um suporte á decisão médica
\section*{Como correr o código}
Para testar os componentes desenvolvidos, os programas encontram-se estruturados da seguinte forma na pasta do projeto:

\begin{itemize}
    \item \textbf{Interface Gráfica (app\_oncologia.m):} Executar o comando \texttt{app\_oncologia} no terminal do MatLab para abrir uma interface gráfica que permite carregar os dados, escolher pacientes e executar o fluxo de trabalho visualmente
    \item \textbf{Demonstração Consola (main\_oncologia.m):} Executar o comando \texttt{main\_oncologia} para rodar o pipeline completo na consola com métricas (Accuracy) e visualização gráfica da Matriz de Confusão.
    \item \textbf{Testes Unitários:} Para validar isoladamente a lógica de cada algoritmo com dados controlados, correr os scripts: \texttt{teste\_bloom.m}, \texttt{teste\_minhash.m} e \texttt{teste\_naivebayes.m}.
\end{itemize}

\textit{*Nota: Os scripts utilizam o dataset real genomic\_variants.csv (ou clinvar.csv). O script preparacao\_dados\_clinvar.m faz o pré-processamento de limpeza, extração de genes/consequências e a partição aleatória dos dados em 80\% para treino e 20\% para teste.}

\section*{Aplicação Conjunta (main\_oncologia.m e app\_oncologia.m)}
O pipeline integrado do sistema junta os 3 módulos para processar sob a seguinte lógica:

\vspace{0.5cm}
\noindent\textbf{Filtro de Bloom (Etapa 1 - Triagem de Risco)}
\begin{itemize}
    \item Verifica se o identificador único da mutação (a chave \texttt{variant\_key} no formato \texttt{CHROM:REF>ALT}) existe na base de dados de treino das variantes classificadas como patogénicas/alto risco.
    \item Permite uma resposta instantânea. Se não estiver no filtro, o médico sabe com 100\% de certeza que a mutação é de baixo risco ou inédita no histórico.
\end{itemize}

\noindent\textbf{MinHash (Etapa 2 - Similaridade de Histórico)}
\begin{itemize}
    \item Converte o perfil textual longo do paciente (que junta gene, cromossoma, bases originais/alteradas e descrição da consequência clínica) em assinaturas compactadas usando 150 hashes lineares.
    \item Estima a similaridade de Jaccard e apresenta na tabela o Top 5 de pacientes do histórico estatisticamente mais parecidos para sugerir tratamentos com base no sucesso de casos anteriores.
\end{itemize}
\newpage
\noindent\textbf{Naïve Bayes (Etapa 3 - Diagnóstico Probabilístico)}
\begin{itemize}
    \item Preve a classe de significância clínica final da mutação (ex: \textit{Benign}, \textit{Pathogenic}, \dots) calculando probabilidades a priori a partir dos atributos genómicos (\texttt{Gene}, \texttt{CHROM}, \texttt{REF}, \texttt{ALT}, \texttt{Consequence}).
    \item Aplica a suavização de Laplace para prevenir que mutações ou genes desconhecidos anulem as frações probabilísticas de classificação.
\end{itemize}

\section*{Testes e Validação}

\subsection*{Módulo 1 - Filtro de Bloom (teste\_bloom.m)}
Para validar o módulo de pertença, inicializámos um filtro de Bloom de tamanho $m = 8000$ e com $k = 3$ funções de hash lineares. Inserimos $n = 1000$ chaves genéticas de treino criadas aleatoriamente.

Efetuámos testes exaustivos realizando $10.000$ consultas com strings que sabíamos de antemão não estarem presentes no filtro para calcular a taxa de falsos positivos de forma prática (empírica).

\vspace{0.5cm}
\noindent\textbf{Resultados e Fórmulas} \\
A probabilidade teórica de falsos positivos ($P_{fp}$) é calculada através da fórmula:

\[
P_{fp} \approx \left(1 - e^{-\frac{k \cdot n}{m}}\right)^k
\]

Para os valores do teste ($m=8000$, $n=1000$, $k=3$):
\[
P_{fp} \approx \left(1 - e^{-\frac{3 \cdot 1000}{8000}}\right)^3 = \left(1 - e^{-0.375}\right)^3 \approx 0.0306 \text{ (ou seja, 3.06\%)}
\]

Nos testes empíricos, os resultados obtidos foram:
\begin{itemize}
    \item \textbf{Falsos Negativos:} 0 (Garantido teoricamente pelo Filtro de Bloom).
    \item \textbf{Taxa Empírica de Falsos Positivos:} $\approx 3.02\%$ (Média observada).
    \item \textbf{Erro de Convergência:} $< 0.1\%$ (Diferença entre o valor teórico e o prático).
\end{itemize}

\begin{figure}[htbp]
  \centering
  \begin{minipage}{0.48\textwidth}
    \centering
    \includegraphics[width=\textwidth]{fig_bloom_fp_vs_k_new.png}
    \caption{Taxa de FP vs. número de hashes $k$.}
    \label{fig:bloom-fp-vs-k}
  \end{minipage}
  \hfill
  \begin{minipage}{0.48\textwidth}
    \centering
    \includegraphics[width=\textwidth]{fig_bloom_fp_vs_m_new.png}
    \caption{Taxa de FP vs. tamanho do filtro $m$.}
    \label{fig:bloom-fp-vs-m}
  \end{minipage}
\end{figure}

\noindent\textbf{Conclusão} \\
Os testes demonstram que a taxa empírica converge de forma muito precisa para a taxa teórica de $3.06\%$. Isto prova que as funções de dispersão (hash) baseadas no algoritmo djb2 distribuem os dados uniformemente pelo intervalo do vetor de bits. Para reduzir as colisões/falsos positivos em ambiente clínico real, o tamanho $m$ do filtro deve ser aumentado.

\subsection*{Módulo 2 - MinHash (teste\_minhash.m)}
Para testar a estimativa de similaridade de Jaccard do MinHash, definimos duas strings representativas de perfis genéticos com uma interseção conhecida de palavras.

Calculámos a similaridade teórica real no papel dividindo o tamanho da interseção dos shingles das duas strings pela sua união:
\[
J(A,B) = \frac{|A \cap B|}{|A \cup B|}
\]

Gerámos depois as assinaturas MinHash para as duas strings usando $N = 200$ hashes lineares e $k_{shingle} = 3$.

\vspace{0.5cm}
\noindent\textbf{Resultados e Conclusão} \\

\begin{figure}[htbp]
  \centering
  \begin{minipage}{0.48\textwidth}
    \centering
    \includegraphics[width=\textwidth]{fig_minhash_error_vs_N_new.png}
    \caption{Erro absoluto médio do MinHash vs. $N$.}
    \label{fig:minhash-error-vs-N}
  \end{minipage}
  \hfill
  \begin{minipage}{0.48\textwidth}
    \centering
    \includegraphics[width=\textwidth]{fig_minhash_k_compare_new.png}
    \caption{Correlação: Jaccard Real vs. Estimado.}
    \label{fig:minhash-k-compare}
  \end{minipage}
\end{figure}
O teste demonstram que a similaridade calculada a partir do vetor de assinaturas compactas do MinHash aproxima-se com elevada precisão do Jaccard real das strings (ex: erro absoluto médio $< 0.02$). O uso de $k_{shingle}=3$ (ou superior) revelou-se ideal para anotações genómicas, pois shingles muito curtos ($k=2$) geram falsas similaridades em palavras curtas do texto clínico.

\subsection*{Módulo 3 - Naïve Bayes (teste\_naivebayes.m)}
Validámos o comportamento probabilístico do classificador de Bayes criando uma tabela pequena com 6 amostras simuladas de 3 genes (\textit{SAMD11}, \textit{EGFR}, \textit{TP53}) e 2 classes (\textit{Pathogenic} e \textit{Benign}).

Treinámos o classificador e testámos a previsão para dois cenários:
\begin{enumerate}
    \item \textbf{Variante comum (Gene e Consequência conhecidos):} O modelo calcula as frequências e classifica a amostra na classe mais provável com elevada certeza probabilística.
    \item \textbf{Variante inédita (Gene desconhecido):} Onde se inseriu o gene \textit{BRCA1} (inexistente no treino).
\end{enumerate}

\vspace{0.5cm}
\noindent\textbf{Resultados e Conclusão} \\
Para a variante inédita, sem a suavização de Laplace, o produto das verosimilhanças seria nulo ($P(\text{Gene}=\text{BRCA1}|C)=0$), provocando uma falha de cálculo. Com o parâmetro $\alpha = 1$, a verosimilhança fallback foi corretamente calculada e o modelo conseguiu classificar com sucesso a amostra baseando-se nas restantes características da mutação (como a Consequência clínica). Isto comprova a robustez e estabilidade do Naïve Bayes para variantes raras.

Além disso, notou-se um grande desequilíbrio nos dados do ClinVar, onde a classe \textit{Likely\_benign} representa a esmagadora maioria dos registos. Esta diferença afeta as decisões do modelo, fazendo com que ele escolha quase sempre a classe com mais exemplos (o que explica a alta taxa de acerto de $87.80\%$ na \textit{Likely\_benign}, enquanto as classes mais raras ficam com $0\%$). Para resolver este problema no futuro, podemos reduzir a quantidade de dados da classe maior (\textit{undersampling}), aumentar os dados das classes mais pequenas (\textit{oversampling}) ou equilibrar as probabilidades iniciais do modelo. Assim, o modelo será mais capaz de detetar as variantes que causam doenças graves.

\begin{figure}[htbp]
  \centering
  \begin{minipage}{0.48\textwidth}
    \centering
    \includegraphics[width=\textwidth]{fig_nb_accuracy_vs_train_new.png}
    \caption{Curva de aprendizagem do Naïve Bayes.}
    \label{fig:nb-accuracy-vs-train}
  \end{minipage}
  \hfill
  \begin{minipage}{0.48\textwidth}
    \centering
    \includegraphics[width=\textwidth]{fig_nb_confusion_new.png}
    \caption{Matriz de confusão no dataset real.}
    \label{fig:nb-confusion}
  \end{minipage}
\end{figure}

\subsection*{Análise de Eficiência Temporal}
Para demonstrar a viabilidade clínica e a escalabilidade do sistema em tempo real, os tempos de execução do treino e da inferência foram medidos com as funções \texttt{tic} e \texttt{toc} do MatLab. O processo completo de treino dos três modelos probabilísticos (sobre 4.000 variantes genómicas) demorou apenas \textbf{0.9665 segundos}. Os tempos médios observados por consulta em fase de inferência estão sumariados na tabela seguinte:

\begin{table}[htbp]
  \centering
  \begin{tabular}{|l|r|}
    \hline
    \textbf{Componente / Algoritmo} & \textbf{Tempo Médio} \\
    \hline
    Treino Total (4.000 registos) & 0.9665 s \\
    Filtro de Bloom (Triagem de Risco) & 0.0420 ms \\
    Naïve Bayes (Predição Diagnóstica) & 1.3372 ms \\
    MinHash (Top 5 Semelhantes em 4.000 do histórico) & 8.8893 ms \\
    \hline
  \end{tabular}
  \caption{Tempos médios de computação das diferentes etapas do pipeline.}
  \label{tab:tempos}
\end{table}

A triagem inicial com o Filtro de Bloom oferece uma resposta extremamente rápida (menos de $0.05$ ms), permitindo filtrar variantes comuns instantaneamente. A inferência probabilística com Naïve Bayes e a busca de histórico com MinHash demoram respetivamente $1.3$ ms e $8.9$ ms, comprovando a elevada eficiência e adequação do sistema para contextos de diagnóstico oncológico em tempo real.

\section*{Conclusão}
Este projeto consolidou a aplicação prática de métodos probabilísticos em problemas complexos da área da saúde. A implementação demonstrou a eficácia da combinação entre algoritmos de eficiência (Bloom e MinHash) e de inferência estatística (Naïve Bayes). A avaliação experimental validou a precisão dos modelos sob condições controladas, destacando a importância da suavização de Laplace e do ajuste de parâmetros de hashing para garantir resultados fidedignos em contextos oncológicos.

\end{document}
